\documentclass[10pt,a4paper]{amsart}
\usepackage{amsmath,amssymb,amsthm,graphicx}
% Retained minimal preamble for the Stacks Project chapter more-morphisms.tex (More on Morphisms).
% Verbatim excerpt of the chapter preamble: Xy-pic setup (native preamble lines 14-18) and the native
% theorem declarations with the mathematical macros (native preamble lines 154-217). Copied mechanically;
% nothing added, removed or redefined. The wrapper class and the amsmath/amssymb/amsthm/graphicx base
% packages are supplied by the assembler, so the class/externaldocument part of the native preamble is omitted.
% For commutative diagrams we use Xy-pic
\usepackage[all]{xy}

% We use 2cell for 2-commutative diagrams.
\xyoption{2cell}
% Theorem environments.
%
\theoremstyle{plain}
\newtheorem{theorem}[subsection]{Theorem}
\newtheorem{proposition}[subsection]{Proposition}
\newtheorem{lemma}[subsection]{Lemma}

\theoremstyle{definition}
\newtheorem{definition}[subsection]{Definition}
\newtheorem{example}[subsection]{Example}
\newtheorem{exercise}[subsection]{Exercise}
\newtheorem{situation}[subsection]{Situation}

\theoremstyle{remark}
\newtheorem{remark}[subsection]{Remark}
\newtheorem{remarks}[subsection]{Remarks}

\numberwithin{equation}{subsection}

% Macros
%
\def\lim{\mathop{\mathrm{lim}}\nolimits}
\def\colim{\mathop{\mathrm{colim}}\nolimits}
\def\Spec{\mathop{\mathrm{Spec}}}
\def\Hom{\mathop{\mathrm{Hom}}\nolimits}
\def\Ext{\mathop{\mathrm{Ext}}\nolimits}
\def\SheafHom{\mathop{\mathcal{H}\!\mathit{om}}\nolimits}
\def\SheafExt{\mathop{\mathcal{E}\!\mathit{xt}}\nolimits}
\def\Sch{\mathit{Sch}}
\def\Mor{\mathop{\mathrm{Mor}}\nolimits}
\def\Ob{\mathop{\mathrm{Ob}}\nolimits}
\def\Sh{\mathop{\mathit{Sh}}\nolimits}
\def\NL{\mathop{N\!L}\nolimits}
\def\CH{\mathop{\mathrm{CH}}\nolimits}
\def\proetale{{pro\text{-}\acute{e}tale}}
\def\etale{{\acute{e}tale}}
\def\QCoh{\mathit{QCoh}}
\def\Ker{\mathop{\mathrm{Ker}}}
\def\Im{\mathop{\mathrm{Im}}}
\def\Coker{\mathop{\mathrm{Coker}}}
\def\Coim{\mathop{\mathrm{Coim}}}

% Boxtimes
%
\DeclareMathSymbol{\boxtimes}{\mathbin}{AMSa}{"02}

%
% Macros for moduli stacks/spaces
%
\def\QCohstack{\mathcal{QC}\!\mathit{oh}}
\def\Cohstack{\mathcal{C}\!\mathit{oh}}
\def\Spacesstack{\mathcal{S}\!\mathit{paces}}
\def\Quotfunctor{\mathrm{Quot}}
\def\Hilbfunctor{\mathrm{Hilb}}
\def\Curvesstack{\mathcal{C}\!\mathit{urves}}
\def\Polarizedstack{\mathcal{P}\!\mathit{olarized}}
\def\Complexesstack{\mathcal{C}\!\mathit{omplexes}}
% \Pic is the operator that assigns to X its picard group, usage \Pic(X)
% \Picardstack_{X/B} denotes the Picard stack of X over B
% \Picardfunctor_{X/B} denotes the Picard functor of X over B
\def\Pic{\mathop{\mathrm{Pic}}\nolimits}
\def\Picardstack{\mathcal{P}\!\mathit{ic}}
\def\Picardfunctor{\mathrm{Pic}}
\def\Deformationcategory{\mathcal{D}\!\mathit{ef}}
\newenvironment{reference}{}{}
\newenvironment{slogan}{}{}
\newenvironment{history}{}{}
\usepackage{cleveref}
\crefname{originallemma}{original lemma}{original lemmas}
\Crefname{originallemma}{Original lemma}{Original lemmas}
\makeatletter
\newlabel{lemma-equivalence-categories-schemes-over-pushout}{{07RV}{source}{}{sourcepointer-1}{}}
\newlabel{lemma-equivalence-categories-schemes-over-pushout@cref}{{[originallemma][0][]07RV}{[1][0][]source}{}{}{}}
\makeatother
\crefname{originallemma}{original lemma}{original lemmas}
\Crefname{originallemma}{Original lemma}{Original lemmas}
\makeatletter
\newlabel{more-algebra-lemma-properties-algebras-over-fibre-product}{{08KQ}{source}{}{sourcepointer-2}{}}
\newlabel{more-algebra-lemma-properties-algebras-over-fibre-product@cref}{{[originallemma][0][]08KQ}{[1][0][]source}{}{}{}}
\makeatother
\title{Flatness and its relatives through the pushout equivalence of schemes over a fibre square}
\author{The Stacks Project Authors}
\date{stacks-project:more-morphisms.tex:local-snapshot-20261011; GNU FDL 1.2 or later}
\begin{document}
\maketitle
\noindent\textit{Editorial note.} Statement span is lines 4402-4423 of sources/stacks-more-morphisms/pinned-original-source.tex; proof span is lines 4425-4430, the complete author proof (4 body lines).\par\medskip
\noindent\textit{Editorial note.} Navigation stubs: 2 cited-but-not-included labels are supplied as non-load-bearing pointers with their pinned tags where the snapshot tag index has them.\par\medskip
\noindent\textit{Editorial note.} extra\_packages is empty on purpose: the retained preamble already loads xy with options and a second declaration clashes in the assembler wrapper.\par\medskip
\noindent\textit{Editorial note.} Text is retained verbatim, including the author's own typos and placeholders.\par\medskip
\noindent\textit{Allowed standard prerequisites (editorial).}\begin{itemize}
\item Explicit prerequisite boundary. The pushout equivalence setup (lemma-equivalence-categories-schemes-over-pushout, same chapter, collected as S3-morph-equivalence-categories-pushout L1424) and the algebra lemma characterizing finite type/flatness/smoothness/etaleness of algebras over fibre-product rings (More on Algebra); cited inputs with the affine-localization reduction written in the retained span.
\end{itemize}\medskip
\ifdefined\hypertarget\hypertarget{sourcepointer-1}{}\fi
\noindent\textit{Source reference.} Original lemma 07RV: sources/stacks-more-morphisms/pinned-original-source.tex, source lines 4277--4293. Cited by the retained statement or proof but not part of the collected spans; supplied only as a navigation pointer. The argument does not depend on its content.\par\medskip
\ifdefined\hypertarget\hypertarget{sourcepointer-2}{}\fi
\noindent\textit{Source reference.} Original lemma 08KQ: sources/stacks-more-morphisms/pinned-original-source.tex, source lines 4430--4430. Cited by the retained statement or proof but not part of the collected spans; supplied only as a navigation pointer. The argument does not depend on its content.\par\medskip
\begin{lemma}
\label{lemma-equivalence-categories-schemes-over-pushout-flat}
In the situation of
Lemma \ref{lemma-equivalence-categories-schemes-over-pushout}.
If $V' = G(V, U', \varphi)$ for some triple $(V, U', \varphi)$, then
\begin{enumerate}
\item $V' \to Y'$ is locally of finite type if and only if $V \to Y$ and
$U' \to X'$ are locally of finite type,
\item $V' \to Y'$ is flat if and only if $V \to Y$ and $U' \to X'$ are flat,
\item $V' \to Y'$ is flat and locally of finite presentation if and only if
$V \to Y$ and $U' \to X'$ are flat and locally of finite presentation,
\item $V' \to Y'$ is smooth if and only if $V \to Y$ and $U' \to X'$ are smooth,
\item $V' \to Y'$ is \'etale if and only if $V \to Y$ and $U' \to X'$
are \'etale, and
\item add more here as needed.
\end{enumerate}
If $W'$ is flat over $Y'$, then the adjunction mapping
$G(F(W')) \to W'$ is an isomorphism. Hence $F$ and $G$ define mutually
quasi-inverse functors between the category of schemes flat over $Y'$
and the category of triples $(V, U', \varphi)$ with $V \to Y$
and $U' \to X'$ flat.
\end{lemma}

\begin{proof}
Looking over affine pieces the assertions of this lemma
are equivalent to the corresponding assertions of
More on Algebra, Lemma
\ref{more-algebra-lemma-properties-algebras-over-fibre-product}.
\end{proof}

\end{document}
